\documentclass[11pt,a4paper]{article}
\usepackage[utf8]{inputenc}
\usepackage[T1]{fontenc}
\usepackage{lmodern}
\usepackage{amsmath,amssymb,bm,array,geometry,color,fancyhdr,hyperref}
\geometry{left=19mm,right=19mm,top=18mm,bottom=19mm}
\definecolor{ink}{RGB}{20,52,72}
\definecolor{accent}{RGB}{22,115,139}
\hypersetup{colorlinks=true,urlcolor=accent,linkcolor=ink,pdfborder={0 0 0}}
\setlength{\parindent}{0pt}\setlength{\parskip}{4pt}
\setlength{\abovedisplayskip}{7pt}\setlength{\belowdisplayskip}{7pt}\setlength{\jot}{4pt}
\setlength{\emergencystretch}{2em}
\pagestyle{fancy}\fancyhf{}\renewcommand{\headrulewidth}{0pt}
\fancyfoot[L]{\small\color{ink}pysnspd / formulación 1.2.1 / 19 de octubre de 2026}
\fancyfoot[R]{\small\thepage}
\newcommand{\block}[1]{\par\vspace{6pt}{\large\bfseries\color{ink}#1}\par\vspace{3pt}}
\newcommand{\critical}[1]{\par{\small\raggedright\textbf{\color{accent}Condición crítica.} #1\par}}
\newcommand{\origin}[1]{\par{\footnotesize\raggedright\color{ink}\textbf{Procedencia:} #1\par}}
\newcommand{\K}[1]{\href{https://eprints.lancs.ac.uk/id/eprint/140252/3/Binder1.pdf\#page=#1}{K20, PDF~#1}}
\newcommand{\A}[1]{\href{https://thesis.caltech.edu/13748/}{A20, p.~#1}}
\newcommand{\V}[1]{\href{https://arxiv.org/abs/1611.06060}{V17, ec.~#1}}
\newcommand{\sref}[2]{\href{Supuestos.pdf\#nameddest=#1}{#1 (Supuestos, p.~#2)}}
\newcommand{\pref}[2]{\href{Parametros.pdf\#nameddest=#1}{#1 (Parámetros, p.~#2)}}
\newcommand{\principalg}{\href{Principal.pdf\#nameddest=G-final}{Principal, cierre de rigidez $G$}}
% Resolver únicamente esta macro al ensamblar la formulación de G:
\newcommand{\gestado}{\textbf{Cierre candidato sujeto a admisibilidad:}\ \principalg; capacidad y rigidez principal positivas, sin pisos artificiales.}
\newcommand{\ptrow}[4]{\hypertarget{#1}{\textbf{#1}}&\( \displaystyle #2 \)& #3 & #4\\[4pt]}
\newenvironment{ptable}{\par\small\setlength{\tabcolsep}{4pt}\renewcommand{\arraystretch}{1.12}
\begin{tabular}{@{}p{.075\linewidth}p{.225\linewidth}p{.285\linewidth}p{.355\linewidth}@{}}\hline
\textbf{ID}&\textbf{Símbolo / unidades}&\textbf{Valor / elección}&\textbf{Función / procedencia}\\\hline}
{\hline\end{tabular}\par}
\newcommand{\assumption}[2]{\par\vspace{6pt}\hypertarget{#1}{\textbf{\color{ink}#1. #2}}\par}
\newcommand{\smallnote}[1]{{\small #1\par}}

\hypersetup{pdftitle={Supuestos - pysnspd 1.2.1},pdfauthor={pysnspd}}
\begin{document}
{\LARGE\bfseries\color{ink}Supuestos y límites del modelo 1.2.1}\par
{\large Registro explícito de hipótesis, dependencias y admisión}\par
\smallnote{Companion de \href{Parametros.pdf}{Parámetros}, \href{Principal.pdf}{Principal}, \href{Auxiliar.pdf}{Auxiliar} y \href{Precomputo.pdf}{Precomputo}. Las entradas experimentales motivan el caso; los cierres propios requieren validación propia. No se presenta una evolución microscópica exacta de Usadel.}
\block{Material y calibración estática}
\assumption{S01}{Resistencia normal inferida de RRR}
\begin{gather*}
 \mathrm{RRR}\ \overset{H}{=}\ \frac{R_{\square,300K}}{R_{\square,N}},
 \qquad R_{\square,N}=\frac{340}{0.8}=425\ \Omega .
\end{gather*}
\smallnote{K20 publica 340 $\Omega$ a temperatura ambiente y RRR 0.8; no define allí la temperatura baja de referencia. Adoptar 425 exige esta convención y tratar el film depositado como representativo de la cinta procesada. Los 608 $\Omega$ de la Tabla suplementaria 1 son una entrada histórica de simulación. Sin medición adicional, no convertirlos en resistencia normal medida ni atribuir incertidumbre cero. \pref{P05}{1}, \pref{P06}{1}; \K{9}--\K{10}, \K{33}.}
\assumption{S02}{Depairing de resonador transferido a la cinta}
\begin{gather*}
 I_{\rm dep}^{120}=38.8\ \mu{\rm A},\qquad
 I_{\rm dep}^{80}=I_{\rm dep}^{120}\frac{80}{120}
 =25.8666667\ \mu{\rm A}.
\end{gather*}
\smallnote{K20 obtiene 38.8 $\mu$A de resonadores fabricados con la misma película. El escalado supone espesor, material y densidad crítica uniformes; ignora un offset de ancho eficaz y constricciones. La referencia se adopta al baño de 0.9 K del caso, sin inventar barras de error. El antecedente Frasca usa 1.05 K y NbN nominal de 6 nm: no demuestra identidad completa de muestra. Depairing y conmutación son observables distintos. \pref{P07}{1}--\pref{P09}{1}; \K{13}; \href{https://arxiv.org/pdf/1904.08596\#page=6}{Frasca, Tabla I}.}
\assumption{S03}{Difusión efectiva calibrada, no medida de transporte}
\smallnote{Se ajusta $D$ para que el máximo homogéneo del cierre y cuadratura nuevos coincida con \pref{P08}{1}. La relación $I_{\rm dep}\propto(R_\square\sqrt D)^{-1}$ permite reescalado sólo si el coeficiente constitutivo DC permanece fijo. Cambiar rigidez, cuadratura o renormalización de $T_c$ exige verificar esa condición o repetir el solver. El ajuste DC de 1.1.3 conserva el funcional estático anterior: se verificaron $D_{425}=1.4103757147386595$ y $D_{608}=0.6891369283$ cm$^2$/s bajo esas hipótesis. \pref{P10}{2}.}
\assumption{S04}{BCS débil y relación de Einstein}
\begin{gather*}
 N_0=(2e^2DR_\square d_f)^{-1},\qquad
 \Delta_{\rm BCS}(0)\simeq1.764k_BT_c .
\end{gather*}
\smallnote{DOS monoespín y material sucio isotrópico efectivo. No variar $D,R_\square,N_0$ de manera independiente ni confundir espesor físico con resolución. La aproximación no determina por sí sola los parámetros microscópicos de NbN fuertemente acoplado. \pref{P12}{2}--\pref{P14}{2}; \href{https://doi.org/10.1103/PhysRev.108.1175}{BCS}; \href{https://arxiv.org/pdf/1904.08596\#page=5}{Frasca, ec. 5}.}
\vfill
\origin{Estos supuestos seleccionan una referencia trazable; la coincidencia estática de una corriente no valida latencia, barrera de detección, cinética ni ancho efectivo.}
\newpage

\block{Geometría y preparación fotónica}
\assumption{S05}{Sección interior de una cinta uniforme}
\smallnote{Cinta de 80 nm, película nominal de 7 nm y longitud activa total de 5 $\mu$m; dominio 2D interior de referencia $L=300$ nm. Sus cortes continúan el mismo superconductor con corriente, no son contactos metálicos ni reservorios BCS sin fase. Mantener $I_{\rm DC}$ físico al comparar longitudes, no la fase total. Absorción longitudinal fija; no añadir jitter geométrico longitudinal. \pref{P03}{1}, \pref{P04}{1}, \pref{P30}{6}; \K{9}, \K{23}.}
\assumption{S06}{Cascada inicial omitida y preparación FD/BE común}
\begin{gather*}
 \Delta_i(0)=\Delta_i^{DC},\qquad
 p_i(0)=f_{\rm FD}(E_i^{DC},T_i^0),\qquad
 n_i(0)=n_{\rm BE}(\Omega,T_i^0).
\end{gather*}
\smallnote{Se comienza después de la conversión electrónica/fonónica inicial. Una temperatura común $T_i^0$ se obtiene al invertir la energía del modelo; el gap inicial sigue siendo el DC. No se afirma equilibrio completo ni coherencia microscópica de esa instantánea. Después evolucionan ocupaciones y condensado, sin imponer FD/BE común. El reloj $t_0=0$ no contiene la duración óptica ni un tiempo de cascada resuelto. \pref{P24}{4}; \A{80}, \A{101}.}
\assumption{S07}{Fracción retenida provisional}
\begin{gather*}
 E_{\rm ret}=\eta_{\rm ret}\frac{h_{\rm P}c}{\lambda},\qquad
 \eta_{\rm ret}=0.8\ {\rm nominal},\quad
 \eta_{\rm ret}\in\{0.667,0.9,1\}\ {\rm alternativas}.
\end{gather*}
\smallnote{0.8 es elección exploratoria, no ajuste ya obtenido de K20. 0.667 reproduce la fracción histórica usada por su simulación; 1 es límite superior de retención en este esquema, no eficiencia óptica medida. La fracción ya retira pérdidas de la cascada omitida: no descontarlas una segunda vez del depósito. El escape dinámico posterior sigue contabilizándose. \pref{P21}{4}, \pref{P22}{4}; \K{33}.}
\assumption{S08}{Anchura gaussiana libre e independiente de $D_{\rm cal}$}
\begin{gather*}
 a_i\propto e^{-|\mathbf r_i-\mathbf r_0|^2/(2\sigma^2)},\qquad
 \sigma=7.549834435\ {\rm nm},\qquad
 \sigma_{\rm dif}(D,\tau)=\sqrt{2D\tau}.
\end{gather*}
\smallnote{La identidad heredada con $D=0.5$ cm$^2$/s y $\tau=0.57$ ps es referencia dimensional. Aquí $\sigma$ permanece parámetro de preparación independiente: no se reescala implícitamente al calibrar $D$ y no es tamaño de foco óptico medido. Normalizar con volúmenes duales y excluir cortes garantiza energía total prescrita; cambiar $\sigma$ cambia la densidad y el reparto e/ph de la preparación común. No sumar 0.57 ps a la latencia. \pref{P23}{4}.}
\vfill
\origin{Las condiciones FD/BE y el perfil gaussiano son hipótesis propias. Referencias A20 distinguen inicialización térmica y conversión inicial; no determinan un depósito espacial universal.}
\newpage

\block{Cinética y material fonónico efectivo}
\assumption{S09}{Tiempos libres de relajación}
\begin{gather*}
 \tau_{ee}(T_c)=5\ {\rm ps},\qquad
 \tau_{ep}(T_c)=24.7\ {\rm ps},\qquad
 \tau_{\rm esc}=20\ {\rm ps}.
\end{gather*}
\smallnote{Los dos primeros son calibradores libres con nomenclatura Allmaras; 4 y 5 ps son controles actuales y K20 usa 6 ps. Los 24.7 ps coinciden con su tabla, sin convertirse en medición del film. Los 20 ps de escape son hipótesis de interfaz; el default histórico de la pipeline es 15 ps y debe sobrescribirse explícitamente. $\tau_{ee}$ afecta BGK y movilidad; $\tau_{ep}$ afecta $\alpha^2F$ y movilidad; variar un tiempo exige conservar dichas dependencias. \pref{P15}{3}--\pref{P17}{3}; \A{95}, \K{33}.}
\assumption{S10}{Aproximación acústica tridimensional hasta Debye}
\begin{gather*}
 g_{ph}(\Omega)=\frac{3\Omega^2}{2\pi^2\hbar^3v_{\rm avg}^3},\qquad
 \alpha^2F(\Omega)=\frac{\hbar\Omega^2}{\pi\tau_0^{e-ph}E_u^3},\\
 v_{\rm avg}=4912\ {\rm m/s},\qquad
 \Omega_D=k_B(664\ {\rm K}),\qquad
 \tau_0^{e-ph}=\frac{720\zeta(5)}{\pi^2}\tau_{ep}(T_c).
\end{gather*}
\smallnote{Es un material exploratorio de película, construido con funciones acústicas; no recupera una DOS DFPT propia ni espectros ópticos exactos de NbN. $N_0\alpha^2F/g_{ph}$ determina la retroacción fonónica conservativa. $v_{\rm avg}$ modifica capacidad y retroacción, además de una escala temporal. $\gamma_{\rm ac}$ se deriva de las mismas capacidades. \pref{P18}{3}--\pref{P20}{3}; \A{81}; \V{3}--\V{4}.}
\assumption{S11}{BGK energético y portador de calor constitutivos}
\smallnote{BGK relaja la forma electrónica hacia FD con el mismo momento energético a gap dado; no fuerza el estado electrónico al baño ni resuelve colisiones e-e microscópicas completas. La escala heredada usa $\tau_{\rm BGK}=\tau_{ee}(T_c)T_c/\max(T_E,T_b)$, donde el piso de baño modifica la tasa y no el objetivo energético. El reparto del calor eléctrico/condensado entre estados electrónicos es una elección constitutiva con restricciones de conservación y soporte. La fórmula computada figura en \href{Auxiliar.pdf\#nameddest=carrier-exacto}{Auxiliar, p.~4}; no se deduce sólo de una temperatura equivalente. La fuente de calor conservada desde 1.1.1 usa $f(1-p)$ al aportar calor y $p(1-f)$ al retirarlo; admite $0\le p\le1$ y $T_E>0$, sin logaritmos en el RHS ni clipping. Conserva el momento energético y exige $P(\sum_k\sigma_kh_k-1/T_E)\ge0$: igualdad sólo en FD. El portador anterior podía enfriar estados vacíos y se descartó por falta de viabilidad continua. BGK conserva energía y no duplica esta fuente.}
\smallnote{Referencia del operador de relajación: \href{https://doi.org/10.1103/PhysRev.94.511}{Bhatnagar, Gross y Krook (1954)}; su realización superconductora energética aquí es propia, con alcance efectivo. \sref{S16}{4}.}
\assumption{S12}{Transporte electrónico local; fonones sin difusión lateral}
\smallnote{Transporte electrónico sobre la malla Delaunay--Voronoi con conductancias y flujo conservativos; coeficiente espectral heredado $D_L=1$ sobre el gap local. Los fonones son poblaciones locales por nodo con escape al baño, sin término de transporte espacial lateral. El escape puede enfriar o calentar según la diferencia con BE del baño y se registra con signo. No se resuelve transporte balístico de fonones, acústica del sustrato ni propagación en interfaces.}
\vfill
\origin{Las ecuaciones cinéticas se inspiran en \V{3}--\V{4}; ensamblado en energías nativas, limitación por disponibilidad y transferencias conservativas son realización propia.}
\newpage

\block{Espectro, temperatura observadora y rigidez térmica}
\assumption{S13}{BCS local sin depairing en la DOS}
\begin{gather*}
 E_{ik}=\sqrt{\xi_k^2+|\Delta_i|^2},\qquad
 \Gamma_{\rm DOS}=0,\qquad 0\le p_{ik}\le1,\quad n_{iq}\ge0 .
\end{gather*}
\smallnote{Espectro BCS con gap local; fase y corriente pertenecen al cierre espacial. No se reconstruye una solución retarded/advanced de Usadel con corriente en cada instante. Un gap finito con ocupaciones calientes no demuestra equilibrio térmico ni espectro microscópico válido. \href{https://doi.org/10.1103/PhysRev.108.1175}{BCS}; \href{Principal.pdf}{Principal}.}
\assumption{S14}{Temperatura equivalente como observador de momento}
\begin{gather*}
 U_{qp,i}=4N_0E_u^2V_i\sum_k w_ke_{ik}p_{ik},\qquad
 U_{qp,i}(\Delta_i,p_i)=
 U_{qp,i}(\Delta_i,f_{\rm FD}(E_i,T_{E,i})) .
\end{gather*}
\smallnote{$w_k$ integra $\xi/E_u$ y $e_{ik}=E_{ik}/E_u$, como en Principal. La inversión a gap dado define $T_E$ sin modificar $p$ ni garantizar FD. Excluye energía fonónica, de circuito y de rigidez. $T_{ph}$ es otro observador energético; ninguno constituye un estado térmico impuesto o calibrador adicional.}
\assumption{S15}{Corrección térmica de depareo anclada al DC}
\begin{gather*}
 \mathcal G_i=\kappa(\rho_i,T_{\rm ref})L_i+
 u_i^*[\varphi(\rho_i,T_{E,i}/T_c,\mathcal Q_i)
       -\varphi(\rho_i,T_{\rm ref}/T_c,\mathcal Q_i)],\\
 \mathcal Q_i=L_i/(u_i^*\rho_i),\quad T_{\rm ref}=T_b\ {\rm congelado},\quad b=T_{\rm ref}/T_c,\\
 U=U_{\rm BCS}+G-\sum_iT_{E,i}G_{T_i},\qquad
 S=S_{qp}-\sum_iG_{T_i},\qquad Q_{{\rm rev},i}=T_{E,i}\dot G_{T_i},\\
 C_{{\rm eff},i}=C_{{\rm FD},i}-T_{E,i}G_{T_iT_i}>0,\quad
 G_{L_i}>0,\quad G_{L_i}+2L_iG_{L_iL_i}>0 .
\end{gather*}
\smallnote{Raíz espectral de Matsubara a gap independiente. El matching propio conserva el DC y recupera \(\kappa(\rho,T_E)L\) a gradiente pequeño; no calibra latencia.}
\begin{gather*}
 H_i^{\mathrm{match}}=\kappa(\rho_i,T_{\rm ref})L_i
 -u_i^*\varphi(\rho_i,b,\mathcal Q_i)\ge0,\qquad
 \lim_{T_{E,i}\to\infty}\mathcal G_i=H_i^{\mathrm{match}} .
\end{gather*}
\smallnote{A \(\rho_i,L_i\) fijos persiste elasticidad independiente de \(T_E\): no se recupera rigidez microscópica nula a temperatura infinita. Este matching no se deriva de Clem--Kogan. No es pequeño: en hotspots auditados aporta 46--71\% de \(\mathcal G_i\) y 67--83\% de \(G_{L_i}\). Es una hipótesis mesoscópica exploratoria, sin admisión cuantitativa de latencia.}
\smallnote{El mapeo mezcla gradientes de amplitud y fase y omite \(\nabla\chi_\nu\); la cinética conserva \(q=0\), \sref{S13}{4}. La corriente no FD no se reconstruye con \(T_E\). \(T_{\rm ref}\) queda congelado; cambiarlo exige revalidar. No es Usadel exacto.}
\critical{La frontera de capacidad del cierre anterior motivó la revisión. Rechazar pérdida de capacidad o ellipticidad: sin pisos ni relajación de tolerancias. Pruebas finitas no demuestran validez global.}
\assumption{S16}{Energía finita y portador electrónico conservativo}
\smallnote{Contabilizar trabajo eléctrico, condensado, cuasipartículas, rigidez, fonones, escape y reservorios con la misma energía usada en la preparación. El calor reversible y la disipación son distintos. El portador electrónico es constitutivo, no una colisión microscópica recuperada; exigir soporte y conservación. Balance energético, estabilidad térmica, producción entrópica y consistencia DC son pruebas independientes. En \(\rho=0,L>0\) el ajuste cancela la singularidad logarítmica; en \(\rho=L=0\) la curvatura es direccional y la convención homogénea sólo precondiciona.}
\vfill
\origin{Definición completa de $G$, derivadas y guardias: \principalg\ y \href{Auxiliar.pdf\#nameddest=G-final}{Auxiliar, p.~1}. \href{https://arxiv.org/html/1207.6421}{Clem--Kogan (2012)}, ecs.~7--14 y apéndice D; extensión anclada y diagnóstico independientes propios.}
\newpage

\block{Electricidad, bordes y lectura}
\assumption{S17}{Cuasineutralidad, sin modo de desequilibrio de carga}
\begin{gather*}
 \mathbf A=0,\qquad\alpha_{ij}=0,\qquad
 h_T\ \hbox{no resuelto},\qquad
 \phi_i\ \hbox{variable algebraica}.
\end{gather*}
\smallnote{El potencial satisface continuidad de corriente con conductancia normal fija y supercorriente del cierre. No se propagan carga capacitiva local, modo de desequilibrio de carga ni campo magnético dinámico. $\alpha_{ij}=0$ no impone fase uniforme: la cinta DC porta gradiente de fase finito. El capacitor del circuito exterior sigue siendo un estado dinámico.}
\assumption{S18}{Cortes superconductores y reservorios explícitos}
\smallnote{Preparar el interior con amplitud uniforme y corriente longitudinal constante dentro del error discretizado, con fases de continuación equipotenciales. En el cierre heredado, la amplitud de los cortes se fija al DC y sus ocupaciones al baño; fases siguen el puerto eléctrico. Su calor y trabajo deben figurar en el balance. No tratar estos cortes como una predicción física de contactos metálicos de K20. Verificar distancia a depósito, longitud de dominio y ausencia de perturbaciones artificiales de borde.}
\smallnote{Mantener malla Delaunay--Voronoi y observar convergencia del problema físico a corriente igual. Condiciones laterales de no flujo; el punto longitudinal de absorción permanece fijo. \pref{P30}{6}.}
\assumption{S19}{Circuito completo y partición inductiva congelada}
\begin{gather*}
 (R_b,L_b,C_c,R_l)=(10\ {\rm k}\Omega,1\ \mu{\rm H},100\ {\rm pF},50\ \Omega),\\
 L_{k,\rm ext}=96\ {\rm nH}
 +(L_{\rm act}-L)\frac{\hbar}{2e}
 \left.\frac{dq_{\rm fis}}{dI}\right|_{DC}.
\end{gather*}
\smallnote{Los cuatro primeros valores son supuestos heredados; K20 sólo informa aquí bias-T resistivo sin sus componentes. Los 96 nH son estimación del inductor añadido, distinta del tramo activo. La derivada es diferencial en el nuevo DC, la partición se calcula una vez y se congela durante el pulso. No incorporar una $L_k(t)$ lumped ni duplicar condensado resuelto. Cambiar $D,R_\square,I_{\rm DC}$ o cierre obliga a recalcular el DC y esta partición. \pref{P25}{5}--\pref{P28}{5}; \K{10}--\K{11}.}
\assumption{S20}{Salida de carga sin transferencia instrumental}
\begin{gather*}
 V_{\rm out}=R_l(I_b-I_s),\qquad V_{\rm dev}=\phi_L-\phi_R .
\end{gather*}
\smallnote{Son señales de puertos distintos. No se añaden filtro, ganancia, ruido, digitalización ni transferencia compleja del amplificador. Los anchos nominales 1.5 GHz del CITLF1 y 6 GHz del osciloscopio no identifican una función de transferencia del modelo; no confundirlos. Un registro determinista sin ensemble no es una IRF ni predice jitter FWHM. \pref{P29}{5}; \K{11}--\K{12}.}
\vfill
\origin{Circuito de tres estados y restricciones eléctricas del modelo propio; datos del montaje K20 conservados únicamente como referencia experimental.}
\newpage

\block{Observación, cobertura numérica y trazabilidad}
\assumption{S21}{Gatillo y ventana de observación declarados}
\smallnote{Primer cruce positivo de $V_{\rm out}-V_{{\rm out},DC}$ por 0.1 mV; confirmar durante 0.5 ps sobre piso de histéresis $0.95V_{\rm th}$. Integrar en 2D hasta cruce + 2 ps; relevo según S25--S32. Sin cruce confirmado a 100 ps, registrar censura, no una latencia inventada. Extraer desde estados aceptados e indicar interpolación. El hotbelt es diagnóstico espacial separado: componente conectada con $|\Delta|\le0.2|\Delta^{DC}|$ que une bordes. \pref{P35}{6}--\pref{P37}{6}.}
\smallnote{La ventana corta cambia sólo el horizonte de observación; no acortar constantes del circuito ni eliminar estados lentos. La latencia postcascada no es latencia óptica absoluta. Comparar 775/1550 bajo la misma regla y reloj; K20 usa desplazamiento del pico ajustado de IRF, por lo que la equivalencia entre ambos observables requiere justificación. \K{12}.}
\assumption{S22}{Modelo discretizado y cobertura espectral finita}
\smallnote{El perfil nominal $\Lambda/E_u=32$, $N_\xi=54$, $N_{ph}=96$, orden de transporte 48, $h=5$ nm y $\Delta t_{\max}=0.02$ ps es propuesta heredada de contraste. No acredita automáticamente el cierre nuevo ni el límite continuo. Refinar espectro, malla, dominio y tiempo sobre energía, corriente, señal y cruce. Recalcular el acoplamiento BCS con la misma cuadratura usada en energía/fuerza para conservar $T_c$. No compensar error de resolución mediante calibración física. Anderson acelera el mapa conjunto de gap/poblaciones/circuito sin cambiar el RHS. Una parada por resolución adaptativa insuficiente no es detección ni censura temporal. Los estados aceptados conservan soporte sin clipping; la limitación por disponibilidad debe conservar cada evento. El camino heredado restringe $I_b,I_s>0$ y requiere fallo controlado si se invierte corriente, sin interpretarlo como ley física universal. \pref{P30}{6}--\pref{P34}{6}.}
\assumption{S23}{Conservación, estabilidad y reproducción son pruebas distintas}
\smallnote{Antes de un transiente: DC nuevo, guardia térmica de rigidez, energía inicial conservada y puertos coherentes. Durante: soporte $0\le p\le1,n\ge0$, finitud, conservación de eventos y balance del sistema completo. Después: convergencia de observables y comparación experimental identificable. Una prueba DC no valida calor finito; las guardias energéticas local/global no promueven validación física del espectro, entropía, rigidez o latencia.}
\smallnote{Los parámetros elegidos aquí son hipótesis que preparan una campaña, no un ajuste ya realizado a 4.2 ps o al récord 2.6 ps. El valor experimental 4.2 $\pm$ 0.4 ps a 21.5 $\mu$A corresponde a latencia relativa 1550 menos 775 nm; los récords de jitter pertenecen a otros observables/casos. \K{5}, \K{22}.}
\assumption{S24}{Identidad reproducible y preservación de resultados}
\smallnote{Guardar hash de fuentes, cierre de rigidez, parámetros físicos y numéricos, cuadratura, malla, solver de depairing, DC, partición inductiva, depósito, gatillo y balance. No reanudar archivos históricos con otra energía o material. Datos pesados, cachés, logs y renders bajo scratch; conservar fuentes y resultados previos. Recursos según afinidad/cuotas actuales, máximo 90\% compartido entre procesos e hilos, usuario \texttt{jdiaz}.}
\block{Fuentes primarias}
\smallnote{\href{https://eprints.lancs.ac.uk/id/eprint/140252/3/Binder1.pdf}{Korzh et al., K20}: artículo aceptado y suplemento, páginas PDF desde 1.
\href{https://thesis.caltech.edu/13748/}{Allmaras, A20}: tesis, páginas impresas (PDF = impresa + 12).
\href{https://arxiv.org/abs/1611.06060}{Vodolazov, V17}: ecs. 3--4.
\href{https://arxiv.org/pdf/1904.08596}{Frasca et al.}: ec. 5, Tabla I.
\href{https://arxiv.org/html/1207.6421}{Clem--Kogan (2012), depareo dirty y energía libre}; \href{https://doi.org/10.1103/PhysRev.108.1175}{BCS (1957)};
\href{https://doi.org/10.1103/PhysRevLett.40.1041}{KWT (1978)};
\href{https://doi.org/10.1103/PhysRev.94.511}{BGK (1954)}.
Referencias inspiran componentes; no validan por atribución el ensamblado efectivo propio.}


\clearpage
\hypertarget{dimensional-1d}{}\hypertarget{hybrid-1d}{}
{\LARGE\bfseries\color{ink}Hipótesis del relevo completo 2D a 1D}\par
\assumption{S25}{Dominio conservado; observación exterior limitada}
\smallnote{Núcleo 2D de 300 nm; después del relevo, 1D de 1000 nm con 350 nm añadidos por lado.
La energía que cruza los reservorios figura en el balance; su evolución
fuera del dominio no se resuelve. El indicador de cortes no es cota rigurosa
del error exterior. Conservación no demuestra suficiencia de longitud.
\pref{P42}{7}, \pref{P48}{7}.}
\assumption{S26}{Restricción espacial, no reemplazo constitutivo}
\smallnote{Para \(t>t_s\) se impone \(\partial_y=0\) en gap complejo,
ocupaciones y potencial. Mismo funcional integrado sobre ancho y espesor:
\(V_j=Wd_f\delta x_j\), caras \(Wd_f\). Se conservan energías nativas,
colisiones, capacidad efectiva, rigidez, fases de cortes y circuito de tres
estados. No se importa la cinética térmica hotbelt de Allmaras.
\href{Principal.pdf\#nameddest=dimensional-1d}{Principal}, \pref{P43}{7}.}
\assumption{S27}{Condición inicial con memoria del evento}
\smallnote{Se remapean \(d/d_{\rm DC}\) complejo y poblaciones nativas.
Las tres variables del circuito se copian exactamente. El exceso de energía respecto del DC de cada
malla se iguala, incluyendo la energía inductiva que abandona el circuito lumped, por una actividad escalar pequeña y acotada, sin clipping
ni sustitución por FD/BE. Es reconciliación de discretización, no calor añadido
durante la evolución. Las colas no observadas se inicializan frías, en la rama estable al \(I_s\) instantáneo, con fases empalmadas y el perfil de borde relativo al DC. Esa inicialización se hace una vez al comenzar la réplica; el relevo hereda su evolución posterior. Es una hipótesis adiabática, no una reconstrucción de su historia. Si excede el margen, se rechaza el relevo.
\pref{P45}{7}; \href{Auxiliar.pdf\#nameddest=dimensional-1d}{Auxiliar}.}
\assumption{S28}{Ausencia resuelta de dinámica transversal}
\smallnote{Error pequeño de reconstrucción del campo complejo, poblaciones y patrón de corriente total respecto de flujo longitudinal uniforme;
ningún contorno válido con enrollamiento positivo o negativo.
No basta vorticidad neta cero. En una región normal no se define fase:
el gap usa coordenadas cartesianas. La reconstrucción mezcla error transversal
y de remapeo. \pref{P46}{7}.}
\origin{Modo transversal constante y volúmenes finitos propios.
Antecedente geométrico:
\href{https://thesis.caltech.edu/13748/}{Allmaras (2020)};
\href{https://arxiv.org/pdf/1805.00130}{Allmaras et al., sec. III.A}.
Diferentes condiciones iniciales, cierre cinético y circuito impiden
identificar ambos modelos.}
\clearpage
\assumption{S29}{Solapamiento independiente y reloj}
\smallnote{\(t_c\): cruce positivo confirmado durante 0.5 ps.
La réplica 1D comienza en el primer estado aceptado que registra el cruce;
su desfase, como máximo un paso 2D, queda registrado.
Hasta \(t_s=t_c+2\) ps el resultado oficial es 2D.
Comparación cada 0.1 ps, incluidos inicio y final.
El núcleo oficial 1D se obtiene reproyectando el último 2D; las colas conservan su estado nativo ya evolucionado en la réplica. Se compara el gap módulo una fase global, porque los cortes de referencia del potencial difieren; no se borran gradientes relativos.
Un fallo conserva datos y declara posible error dimensional.
\pref{P44}{7}, \pref{P47}{7}.}
\assumption{S30}{Indicador de truncamiento longitudinal}
\smallnote{Se integra \(|P_L|+|P_R|\) en puntos medios aceptados sin
cancelación entre cortes. Normalización:
\(E_{\rm ret}+\int|I_sV_{\rm dev}|\,{\rm d}t\).
Cociente máximo 0.02 en ambas fases.
Margen práctico inicial, no propiedad universal de NbN ni prueba de exterior
frío. \pref{P48}{7}.}
\assumption{S31}{Aceleración con aceptación física}
\smallnote{El 2D conserva punto medio implícito y guardias. En 1D se compara
un paso completo con dos medios pasos y se acepta la segunda solución sólo
si pasan error escalado y balances independientes.
No se suprimen estados lentos ni se fuerza equilibrio.
Las celdas de borde inválidas se reconstruyen por intersección convexa de Voronoi, sin cambiar sitios ni triangulación. Se conservan capacidad y rigidez positivas, soporte de Pauli/Bose,
tolerancia no lineal y prohibición vigente de inversión de corriente.
\pref{P49}{7}.}
\assumption{S32}{Alcance temporal y validación pendiente}
\smallnote{15 ns es tiempo final solicitado, separado del techo de 100 ps
sin detección. No certifica reset: el circuito puede requerir más tiempo
y el dominio fallar antes. Dos ps de concordancia no demuestran estabilidad
de todos los modos transversales durante 15 ns; su ausencia posterior es
hipótesis explícita. Pruebas cortas verifican operadores, proyección, reinicio
y balance. La campaña física debe validar admisibilidad.
El cambio no valida por sí mismo jitter ni latencia experimental.
\pref{P44}{7}, \pref{P46}{7}--\pref{P49}{7}.}
\origin{Reset y realimentación electro-térmica:
\href{https://arxiv.org/abs/1008.0895}{Annunziata et al. (2010)}.
El cierre efectivo previo y sus límites siguen vigentes:
\href{https://arxiv.org/pdf/1611.06060}{Vodolazov (2017)}
es referencia, no identificación con una teoría exacta.}

\clearpage
\assumption{S33}{Cambio de representación de energía e inductancia}
\smallnote{Los 700 nm añadidos dejan de estar representados por inductancia
lumped y pasan al funcional 1D. Se conserva el inductor añadido de 96 nH.
Las dos particiones se fijan con el mismo DC, no se ajustan dinámicamente:
\(\Delta L=L'_{\rm DC}(L_1-L_2)\).
La película 1D recibe el exceso
\(\delta E_{\rm film,2}+\frac12\Delta L(I_s^2-I_{\rm DC}^2)\);
las tres variables del circuito se copian.
El cambio constante de referencia DC incluye película y circuito.
No se crea calor para compensar una inductancia contada dos veces.
\href{Auxiliar.pdf\#nameddest=dimensional-1d}{Auxiliar};
\pref{P43}{7}.}
\assumption{S34}{Dominio finito y diagnóstico de latching}
\smallnote{1 \(\mu\)m es punto de partida, no longitud suficiente garantizada
para 15 ns. Se mantienen el balance independiente y el indicador de cortes;
se vigilan 25 nm de nodos libres junto a ambos extremos 1D.
Un frente térmico o de supresión persistente termina con posible error de
dominio. Desde 5 ns, un segmento resistivo de al menos 20 nm con plateau
eléctrico y térmico durante 1 ns se guarda como posible latching.
No identifica de forma exclusiva una fase microscópica ni demuestra una avería
numérica. \pref{P50}{8}, \pref{P53}{8};
\href{https://arxiv.org/abs/1008.0895}{Annunziata et al.}.}
\assumption{S35}{Polaridad, límites eléctricos y recuperación}
\smallnote{Se vigilan corrientes globales positivas y techos declarados de
corriente/voltaje; ante incumplimiento se guardan estado y motivo.
Una lectura negativa por el condensador no implica que la corriente global
haya invertido su sentido: no se rechaza automáticamente.
La recuperación anticipada exige un intervalo estable de 1 ns con corriente,
gap, ocupaciones y las tres variables circuitales dentro de márgenes.
15 ns permanece como límite temporal superior, no como recuperación garantizada.
\pref{P51}{8}--\pref{P56}{8};
\href{https://arxiv.org/pdf/1302.2852}{Kerman et al. (2013), fig. 1 y circuito AC}.}
\end{document}
